\documentclass[10pt,a4paper]{amsart}
\usepackage[margin=19mm]{geometry}
\usepackage{amsmath,amssymb,mathtools}
\usepackage[scr=rsfs,cal=euler]{mathalfa}
\usepackage{xfrac}
\usepackage[unicode,hidelinks,bookmarks=false]{hyperref}
\newcommand{\ie}{i.e.\ }
\newcommand{\cf}{cf.\ }
\newcommand{\resp}{respectively\ }
\newcommand{\Subsection}{\subsection}
\providecommand{\nobreakdash}{\nobreak\hbox{-}}
\setlength{\emergencystretch}{2em}
\multlinegap0pt
\usepackage{amssymb}
\usepackage{xcolor}
\newtheorem{lem}{Lemma}
\newcommand{\PHF}[2]{P_{{#2},{#1}}}
\newcommand{\RR}{\mathbb{R}}
\newcommand{\SosHF}[2]{\Sigma_{{#2},{#1}}}
\title{An efficient sum of squares nonnegativity certificate for quaternary quartic}
\author{Dmitrii V. Pasechnik}
\date{Source: arXiv:1511.03473v4 (CC BY 4.0)}
\begin{document}
\maketitle

\noindent\emph{Source and licence.} An efficient sum of squares nonnegativity certificate for quaternary quartic. Version arXiv:1511.03473v4 (CC BY 4.0), distributed by arXiv under the Creative Commons Attribution 4.0 International licence, http://creativecommons.org/licenses/by/4.0/, as recorded in the arXiv licence metadata for this exact version and shown on the abstract page https://arxiv.org/abs/1511.03473v4. Full licence notice: \texttt{licenses/arxiv-1511.03473-CC-BY-4.0.txt}.\par\smallskip

\noindent\textit{Editorial scope.} Counted unit: the article's lem lem1 (source0.tex lines 259--276). The article's complete original proof of it is delivered with it (source0.tex lines 278--294, one \texttt{proof} environment of 17 lines). No mathematics is written, repaired or re-derived: the statement and the proof are the article's own lines, in the article's own notation, and no retained line is substituted or re-flowed.  No further source-local context is needed: every cross-reference of every retained line resolves inside the two delivered spans.  Nothing was omitted from any retained span, so the package carries no omission record.  No retained line contains a citation, so the wrapper carries no bibliography and no bibliography entry.  No retained line contains a figure environment, an image-inclusion command or drawing code, so no image is delivered and the package declares no asset dependency.  Editorial formatting only, in addition: an AMS \texttt{amsart} preamble that restates the article's own declarations and macros used by the retained lines, namely the environments \texttt{lem} on separate counters, together with the standard packages those lines need.  The retained environments print in this document as Lemma 1 in order of appearance, whereas the article prints the same items with the numbering of its own full document.  The complete native source of the article as distributed in the arXiv e-print for this exact version is bundled unchanged as \texttt{sources/100355/source0.tex}.
\begin{lem}\label{lem1}
Let $f\in\PHF{2d}{n}$, 
and $S:=\{v\in \RR^{n}\mid \sum_{k=1}^n v_k^2=1\}$ be the unit sphere in $\RR^{n}$. Denote $u:=x_1$.
\begin{itemize}
\item[(i)] Let $f$ have a zero on $S$.  An orthogonal change
of coordinates bringing the zero to $(1{:}0{:}\dots {:} 0)$ transforms $f$ into the form
\begin{equation}
f(u,x)=\sum_{k=2}^{2d} f_k(x) u^{2d-k}, \quad f_k\in\RR[x_2,\dots,x_n].\label{red0}
\end{equation}
\item[(ii)] Let $f>0$, with a  minimum on $S$
reached at $x^*\in S$. 
An orthogonal change of coordinates bringing $x^*$ to $(1{:}0{:}\dots {:}0)$ transforms $f$ into the form
\begin{equation}
    f(u,x)=u^{2d}+\sum_{k=2}^{2d} f_k(x) u^{2d-k}, \quad f_k\in\RR[x_2,\dots,_n],\text{ so that $f(u,x)-u^{2d}\geq 0$.}\label{rednon0}
\end{equation}
\item[(iii)] $f_2\geq 0$ and $f_{2d}\geq 0$, and moreover, $f_2\in \SosHF{2}{n}$.
\end{itemize}
\end{lem}

\begin{proof}
(i). After the transformation, $f$ cannot have a term $u^{2d}$, as otherwise
it cannot vanish on $(1{:}0{:}\dots {:}0)$. It cannot have a term $f_1(x)u^{2d-1}$, as 
it is nonnegative. Thus we have the required form.

(ii).  Without loss of generality, after the transformation we have
\[
f(u,x)=u^{2d}+g(u,x), \quad g(u,x):=\sum_{k=1}^{2d} f_k(x) u^{2d-k}.
\]
Note that on $S$ one has $f(u,x)\geq 1$, as the minimum of $f$ on $S$ equals 1.
Thus on $S$ one has $g(u,x)=f(u,x)-u^{2d}\geq f(u,x)-1\geq 0$. Hence $g(u,x)\geq 0$,
as it is homogeneous. Moreover, it has a zero, $(1{:}0\dots {:}0)$, on 
$S$, and this (i) applies, ensuring $f_1=0$, as required.

(iii). $f_2\geq 0$ and $f_{2d}\geq 0$ follows from $f\geq 0$ in the case (i), and
from $f(u,x)-u^{2d}\geq 0$ in the case (ii). As $f_2$ is quadratic, it is an s.o.s.
\end{proof}

\end{document}
